% Modernised reading — fonds Alexandre Grothendieck, shelfmark 71
% (pages 1-47, the whole folder).
%
% This is an INTERPRETATION, not a transcription. It re-reads the pages in
% today's notation and vocabulary, states the hypotheses the manuscript leaves
% implicit, and drops the critical apparatus so that the mathematics reads
% straight through. Everything the brackets and underlines carried moves into
% a footnote. It is held to being mathematically correct as it stands; where
% the manuscript is loose or mistaken, the true statement is what appears here
% and the footnote says what the page has. In any dispute of substance the
% transcription governs: whoever cites Grothendieck must cite batch-01.fr.tex
% (pp. 1-20), batch-02.fr.tex (pp. 21-40) or batch-03.fr.tex (pp. 41-47).
%
% Pass: Opus 5.5 (claude-opus-5-5), 2026-10-10 — first pass, unchecked against the pages by a human.
% The folder was read whole, its three batches together; this is one document
% for the shelfmark. It was made on Opus 5.5 and has not been measured against
% the reference reading of folder 115 (made on Opus 5). The literature named
% in the footnotes is cited from memory and was not checked against the
% sources.
%
% Scope. Pages 1, 4, 15, 17, 19 are blank separators; 3 and 14 typed USTL
% « avis de soutenance » (19 avril 1977); 8 a printed bank leaflet (janvier
% 1978); 25, 34, 40, 45 versos of the mémoire (false starts, two pencil
% drawings on 45). None carries mathematics of his; none is described beyond
% the spine section. Pages 20, 22, 24, 26, 28, 30, 32, 33, 35-39, 41-44, 46,
% 47 are a mémoire in another hand (signed, « 15 juillet 1976, Montpellier »).
% Following RIGHTS.md and the transcription's own rule (folder 100's), the
% mémoire is NOT restated beyond what the transcription's notes say: its
% statements are named as the notes name them, its proofs are not given, its
% author is not named. What this reading comments on is his marks and the
% mathematics of the statements.
%
% Spine. Three strands of his own, around the Jordan curve theorem, and his
% reading of the mémoire:
%   pp. 2, 13, 16  theta-graphs: three arcs joining x, y; three regions;
%                  a) <=> a') <=> b) <=> c) (page 13), c) => a) proved,
%                  a) => b) begun; page 16 the same with an index set, and
%                  three equivalent forms of « C separates the surface E ».
%   pp. 9, 11, 7, 5  his numbered sequence 1)-6) (2) absent): complementary
%                  open pairs (regular open sets); a closed K with two
%                  relatively open pieces; pseudo-partitions of X into three
%                  open sets and three walls; X cut open along the inner walls
%                  and reglued (the space 𝒳 -> X); an « immersion lemma »
%                  announced.
%   pp. 6, 12      a cancelled computation: a filtration of a realised poset
%                  by an ordered set I, spaces Σ_d indexed by flags, and the
%                  combinatorial cutting Σ(K, L) of a simplicial complex.
%   p. 10          unrelated: an automorphism t of a group on λ_0, λ_1, λ_2.
%   p. 18          two J-polygons meeting along an arc: exterior or nested; the
%                  glued J-polygon Γ'' and its interior and exterior.
%   pp. 20-47      the mémoire (Théorème 1, polygonal Jordan; Théorème 2,
%                  three good vertices), with his notes on separate leaves
%                  (21, 29, 31) and his marks in the margins.
%
% Notation fixed for the document.
%  - « J-courbe » (ours) for a simple closed curve satisfying the conclusion of
%    the Jordan theorem in the strong local form; the pages say « courbe de
%    Jordan », which today means any simple closed curve. « J-polygone » is
%    the mémoire's word (Définition 1) and is kept.
%  - Frontier of an open set U: U̇ = Ū \ U, his dotted letter. Frontier of a
%    closed K: K̇ = K \ int K. Int, Ext of a J-polygon as on the pages.
%  - Theta-graph: arcs Γ'_1, Γ'_2, Γ'_3, ∂ = {x, y}, Γ''_i = Γ'_i \ ∂ (page 2;
%    page 16 writes Γ°_i), Γ_k = Γ'_i ∪ Γ'_j ({i,j,k} = {1,2,3}), regions
%    V_k, sectors v_k at x, U_k = V_i ∪ V_j ∪ Γ''_k.
%  - Pages 7 and 5: X_i := W ∪ Ū_i (the gluing piece) is kept apart from the
%    complement X \ U_i, written out; the page calls both X_i.
%
% Substantive departures, each footnoted where it occurs:
%  p. 2    « Γ_1 = Γ_2 ∪ Γ_3 » and the last member Γ''_1 ∪ Γ''_2 ∪ ∂ of the
%          first line read Γ_1 = Γ'_2 ∪ Γ'_3 = Γ''_2 ∪ Γ''_3 ∪ ∂.
%  p. 9    1) d) is not equivalent to a)-c) without « U ∪ V dense »
%          (counterexample ours). 3): the conclusions need B ⊂ Ū (true when
%          K is the closure of its interior); counterexample ours (disc with a
%          hair). « relativement fermé dans X » read: B closed in X iff
%          B̄ ∩ C = ∅, as the Cor. then says.
%  p. 11   « pseudo-partition » is not defined on the page; definition ours,
%          from the relations the page lists; a) <=> b) checked under it; the
%          converse « Je reviens... » made explicit (ours).
%  p. 7    X̄_i = X_i ∪ B'_i corrected to (X \ U_i) ∪ B_i; « U_i \ C_i » read
%          U_i \ B_i = Int U_i; the hypothesis B_i ⊂ closure(U_i \ B_i)
%          supplied; the two uses of X_i separated.
%  p. 5    the check asked (« vérifier si ... C_i ∩ K̇ = ∅ ») answered: it is
%          implied by K̇ ⊂ ∪ U_i (ours). The étale / closed-embedding claim
%          stated as checked; the immersion lemma 6) breaks off.
%  p. 6    « φ(M(e)) ≤ i » read max φ(M(e)) ≤ i, as page 12's μ says. Σ_d is
%          never defined; the bijection D(d) ≅ E(d) is reported, not checked.
%  p. 10   the computation headed t^3 computes t∘t; with λ_0λ_1λ_2 = 1
%          (inferred, not written) t² = int(λ_1λ_2) = int(λ_0^{-1}).
%  p. 13   « courbe de Jordan » = J-courbe; arcs supposed polygonal (tame at
%          x); a) => b) supplied from the formulas of pages 2 and 16.
%  p. 16   « U ⊂ C », « V = U - Ū », « Ū ≠ E \ C », « V_1 = U_1 ∩ U_2 » and
%          the U_j, V_j under ∪, ∩ read U ⊂ E, V = E - Ū, U ≠ E \ C,
%          V_1 = U_2 ∩ U_3, V_j, U_j; the « y » of c) read x''.
%  p. 18   Γ ≠ Γ' needed (else Γ'' is not a polygon); the dichotomy proved
%          (ours).
%  p. 29   (i): « Γ ∩ T̊ = [s,s'] ∪ [s,s''] » read Γ ∩ T; the « t ≠ s » of
%          Cor Main read « [s,t] a diagonal »; vertices on ]s',s''[ added to
%          make Cor Main hold when Φ = ∅ (ours).
%  p. 43   the transcription's note glosses § 4 « Preuve du théorème 2 » as
%          « tout polygone est un J-polygone », which is Théorème 1; the
%          reading follows page 20 and the batch-3 header.
%
% Links the pages do not write (ours, and said so in the body): the theta
% lemma of page 13, applied to Γ ∪ [s,s'] for a diagonal [s,s'], is the
% inductive step of the polygonal Jordan theorem, with no condition on the
% diagonal; page 31 provides an exterior diagonal for every non-convex
% polygon without the J- hypothesis (his « (J-) » on Lemme VIII, « inutile »);
% page 18's cases a) and b) are the interior and exterior diagonal of the
% mémoire's §§ 3-4; the pseudo-partition of page 11 is what the theta
% configuration becomes once ∂ is removed; the X cut open of pages 7 and 5 and
% the Σ(K, L) of page 12 are one operation, topological and combinatorial;
% folder 81's Version I (pp. 13-24) is the same theme of « découpage »;
% folder 89's combinatorial polygons (pp. 19-20) are the mémoire's « polygone
% abstrait ».
%
% Announced and never established: a) => b) of page 13 (« Prouvons que
% Γ_3 ... », stops); the meaning of the index set on page 16; the end of
% page 7; the immersion lemma 6) of page 5; what Σ_d is and the bijection
% α_d (pages 6 and 12, cancelled); the question « |K̃| ≅? Σ(|K|,|L|) »
% (page 12, struck); cases (ii) of page 29's lemma.
\documentclass[11pt,a4paper]{article}
\input{../preamble/grothendieck}

\folder{71}
\pages{1}{47}
\dating{[à partir de 1976]}
\watermark{Édition de démonstration\\interprétation personnelle de l'œuvre}
\foldertitle{Théorème de Jordan : notes manuscrites (1976, s.d.) — lecture modernisée du dossier entier}

\begin{document}

\section*{Résumé}

\begin{resume}
Une ficelle fermée posée sur une table, sans se croiser elle-même, sépare la
table en deux : un dedans et un dehors. C'est le théorème de Jordan, l'un des
énoncés les plus évidents de la géométrie et l'un des plus difficiles à
démontrer proprement. Le dossier en garde deux approches. L'une est de sa
main, éparse, sur une douzaine de feuillets ; l'autre est un mémoire soigné
d'une autre main, daté de Montpellier le 15 juillet 1976, qui démontre le
théorème pour les polygones, et qu'il a lu plume en main.

Sa propre approche commence par une figure plus riche que la ficelle : deux
points reliés par trois chemins qui ne se croisent pas, comme la lettre grecque
thêta. Trois chemins font trois boucles, en en prenant deux à la fois, et trois
régions du plan. Il montre que si deux des trois boucles séparent le plan
comme il faut, la troisième aussi. C'est un mécanisme de récurrence : un
polygone qu'on coupe par une diagonale devient un thêta dont les deux autres
boucles ont moins de sommets, et le théorème se transmet de proche en proche.
Les pages ne le disent pas en ces termes ; c'est ce que leurs pièces, mises
bout à bout, permettent de dire.

À côté, il prépare un cadre général, de la topologie la plus élémentaire : deux
ouverts dont chacun est exactement l'extérieur de l'autre, trois régions qui se
touchent deux à deux le long de trois murs, et l'opération qui consiste à
découper un espace le long de murs intérieurs puis à recoller les morceaux.
C'est la langue du découpage, qu'on retrouve ailleurs dans ce groupe de
dossiers (le dossier 81 découpe des espaces par des coupures). Un calcul
abandonné fait la même chose pour des complexes simpliciaux, sommet par
sommet.

Le mémoire, enfin, est l'occasion de voir Grothendieck lecteur. Ses marges
disent « style », « vague », « jargon ? », « C'est une salade… », mais aussi,
précisément, ce qui manque : un « non vides ! » dans une définition, un
« faux pour $n = 3$ ! » sous un lemme, la remarque que c'est la \emph{paire}
de deux arcs et non le \emph{couple} qui est déterminée. Sur des feuillets
séparés, il refait à sa façon deux étapes du mémoire, en enlevant une
hypothèse inutile. On y voit une exigence qui ne tient pas à la difficulté de
l'énoncé : le théorème est élémentaire, et c'est justement pour cela que chaque
hypothèse doit être à sa place.

Les noms modernes à chercher sont ceux du théorème de Jordan polygonal, du
théorème du thêta, des ouverts réguliers et du recollement, et, pour le second
théorème du mémoire, ceux des « oreilles » d'un polygone.

\keywords{Jordan curve theorem, polygonal Jordan theorem, theta-curve theorem, regular open set, locally closed set, cutting and gluing, local homeomorphism, polygon diagonal, convex hull, two ears theorem, Hurwitz action}
\end{resume}

\pagerange{2}{47}
\section*{Le fil du dossier, et les conventions}

\subsection*{Les feuillets}

Le dossier compte quarante-sept pages, sans pagination de sa main. Les pages
1, 4, 15, 17 et 19 sont des intercalaires blancs ; les pages 3 et 14 des avis
de soutenance dactylographiés de l'université de Montpellier (19 avril 1977),
la page 8 un dépliant bancaire imprimé (janvier 1978), sans rien de sa main :
ces versos s'accordent avec la datation « à partir de 1976 » et n'en disent
pas plus. Les feuillets de sa main sont les pages 2, 5 à 7, 9 à 13, 16, 18,
21, 23, 27, 29 et 31. Les pages 20 à 47 sont, pour le reste, un mémoire d'une
autre main, paginé de I à XIX ; ses versos 25, 34, 40 et 45 ne portent que des
faux départs et, page 45, deux croquis au crayon.

\subsection*{Les stations}

\begin{itemize}
  \item \textbf{Pages 2, 13 et 16 : le thêta.} Trois arcs joignant deux points,
    trois courbes fermées, trois régions ; l'équivalence de la page 13 et sa
    démonstration partielle ; page 16, le même dispositif indexé par un
    ensemble quelconque, et trois formes équivalentes de « $C$ sépare la
    surface ».
  \item \textbf{Pages 9, 11, 7, 5 : une suite numérotée 1) à 6)}, dans cet
    ordre (le 2) manque). Ouverts complémentaires ; un fermé $K$ partagé en
    deux pièces ; pseudo-partitions ; découpage le long de murs intérieurs et
    recollement ; un « lemme d'immersion » annoncé.
  \item \textbf{Pages 6 et 12 : un calcul barré}, sur une filtration d'un
    ensemble ordonné réalisé et sur le découpage d'un complexe simplicial le
    long d'un sous-complexe.
  \item \textbf{Page 10 : une page d'algèbre} sans rapport avec le reste.
  \item \textbf{Page 18 : recoller deux J-polygones} le long d'un arc commun.
  \item \textbf{Pages 20 à 47 : le mémoire}, avec ses notes à lui sur des
    feuillets séparés (21, 29, 31) et ses marques en marge.
\end{itemize}

\subsection*{Conventions}

$E$ désigne le plan $\mathbb{R}^2$ (page 16 : une surface connexe). Pour un
ouvert $U$ d'un espace $X$, $\dot{U} = \overline{U} \smallsetminus U$ est sa
frontière, la lettre pointée des pages ; pour un fermé $K$, $\dot{K} = K
\smallsetminus \operatorname{int} K$. Le complémentaire est noté
$\complement$.

Les pages disent « courbe de Jordan » pour une courbe fermée simple
\emph{qui vérifie la conclusion du théorème de Jordan} ; aujourd'hui
l'expression désigne n'importe quelle courbe fermée simple. On écrit donc
\emph{J-courbe} : une courbe fermée simple $C \subset E$ est une J-courbe si $E
\smallsetminus C$ a exactement deux composantes connexes et si, en chaque point
de $C$, les deux côtés locaux de $C$ sont dans des composantes
différentes\footnote{Le terme est de cette lecture ; il est calqué sur le
« J-polygone » de la définition 1 du mémoire (page 22), qui demande exactement
cela, plus l'assertion qu'une composante est bornée et l'autre non — ce qui,
dans le plan, est automatique dès qu'il y a deux composantes, le complémentaire
d'un grand disque étant connexe.}. Le mémoire dit \emph{J-polygone}, et l'on
garde son mot.

Pour le thêta, on fixe : deux points $x \neq y$, $\partial = \{x, y\}$ ; trois
arcs simples $\Gamma'_1, \Gamma'_2, \Gamma'_3$ de $x$ à $y$, deux à deux
disjoints hors de $\partial$ ; leurs intérieurs $\Gamma''_i = \Gamma'_i
\smallsetminus \partial$ ; pour $\{i, j, k\} = \{1, 2, 3\}$, la courbe fermée
$\Gamma_k = \Gamma'_i \cup \Gamma'_j$ ; $\Gamma = \Gamma'_1 \cup \Gamma'_2 \cup
\Gamma'_3$ ; en $x$, le secteur $v_k$ compris entre $\Gamma'_i$ et
$\Gamma'_j$, et, quand elles existent, les régions $V_k \supset v_k$ de $E
\smallsetminus \Gamma$.

\subsection*{Ce que seule une lecture d'ensemble peut dire}

Les pièces de sa main s'emboîtent en une démonstration du théorème de Jordan
pour les polygones, qu'aucune page n'écrit d'un tenant. Si $[s, s']$ est une
\emph{diagonale} d'un polygone $\Gamma$ — un segment joignant deux sommets
non adjacents et ne rencontrant $\Gamma$ qu'en ses extrémités, qu'il passe à
l'intérieur ou à l'extérieur —, alors $\Gamma \cup [s, s']$ est un thêta, dont
les deux autres courbes fermées $\Lambda^1, \Lambda^2$ ont moins de sommets que
$\Gamma$. Par la page 13, si $\Lambda^1$ et $\Lambda^2$ sont des J-courbes,
$\Gamma$ en est une. La page 31 fournit une diagonale \emph{extérieure} à tout
polygone non convexe, sans rien supposer de lui ; un polygone convexe a des
diagonales évidentes ; le triangle est le cas de départ. La page 18 dit ensuite
qui est l'intérieur de $\Gamma$, selon que la diagonale passe dehors (cas b) :
$\Lambda^1, \Lambda^2$ emboîtés) ou dedans (cas a) : extérieurs l'un à
l'autre). C'est, dans un autre ordre, la structure du mémoire, dont le § 3
coupe selon une diagonale extérieure et le § 4 selon une diagonale intérieure ;
et c'est ce qui donne sens aux deux marques de la page 32, où il met entre
parenthèses le « J- » de l'énoncé du lemme VIII et qualifie ce lemme
d'« inutile ». Le rapprochement est de cette lecture.

Deux autres liens, de cette lecture aussi. La pseudo-partition de la page 11
— trois ouverts deux à deux adjacents le long de trois murs sans point triple —
est exactement ce que devient le thêta quand on lui retire ses deux sommets.
Et l'espace $\mathcal{X}$ des pages 7 et 5, obtenu en découpant $X$ le long de
murs intérieurs, est la version topologique de l'opération combinatoire
$\Sigma(K, L)$ de la page 12.

\subsection*{Ce que le dossier annonce sans l'établir}

L'implication a) $\Rightarrow$ b) de la page 13, commencée (« Prouvons que
$\Gamma_3$… ») ; le rôle de l'ensemble d'indices de la page 16 ; la fin de la
page 7 ; le « lemme d'immersion » 6) de la page 5 ; ce que sont les espaces
$\Sigma_d$ et la bijection $\alpha_d$ des pages 6 et 12, barrées ; la question
« $|\tilde{K}| \simeq \Sigma(|K|, |L|)$ ? » de la page 12, barrée ; le cas (ii)
du lemme de la page 29.

\pagerange{2}{16}
\section*{I. Le thêta (pages 2, 13 et 16)}

\pagerange{2}{2}
\subsection*{Trois arcs, trois courbes, trois régions (page 2)}

La page 2 dresse l'inventaire du thêta. Chaque région $V_k$ est bordée par les
deux arcs $\Gamma'_i, \Gamma'_j$ qui l'enserrent, et chaque arc ouvert
$\Gamma''_i$ borde les deux régions $V_j, V_k$ :
\[
\Gamma_k = \Gamma'_i \cup \Gamma'_j = \Gamma''_i \cup \Gamma''_j \cup \partial,
\qquad
U_k = V_i \cup V_j \cup \Gamma''_k, \qquad \dot{U}_k = \Gamma_k ,
\]
pour $\{i, j, k\} = \{1, 2, 3\}$\footnote{La page écrit d'abord $\Gamma_1 =
\Gamma_2 \cup \Gamma_3$, sans accents, ce qui n'a pas de sens pour des courbes
fermées ; la page 13 écrit $\Gamma_1 = \Gamma'_2 \cup \Gamma'_3$, qu'on suit.
Le dernier membre de la première ligne, $\Gamma''_1 \cup \Gamma''_2 \cup
\partial$, d'indices de lecture incertaine, doit être $\Gamma''_2 \cup
\Gamma''_3 \cup \partial$ ; la troisième ligne, et le $\Gamma^{*}_3 =
\Gamma''_1 \cup \Gamma''_2 \cup \partial$ qui suit, sont justes. Les lettres
$\Gamma''_i$ de la page 2 sont les $\Gamma^{\circ}_i$ de la page 16.}. Ces
formules disent ce qu'on attend : $E \smallsetminus \Gamma_k$ a deux
composantes, $V_k$ d'un côté, $U_k$ de l'autre, cette dernière formée des deux
autres régions soudées le long de l'arc $\Gamma''_k$ qui les sépare. Un croquis
de deux cercles sécants et les lettres « $U, C, V$ » rappellent la situation
de la courbe seule.

\pagerange{13}{13}
\subsection*{L'équivalence de la page 13}

On suppose les arcs \emph{polygonaux}\footnote{La page ne dit rien de la
nature des arcs. L'énoncé a) parle de trois secteurs en $x$, ce qui demande
qu'un petit disque centré en $x$ soit coupé par $\Gamma$ en trois secteurs, et
la démonstration fait passer un côté d'un arc d'un point à un autre le long de
l'arc : les deux sont immédiats pour des arcs polygonaux, comme ceux du
mémoire. Pour des arcs quelconques du plan, ils le sont encore, mais par le
théorème de Schoenflies, qui est plus fort que ce qu'on cherche à démontrer.}.
Les conditions suivantes sont équivalentes :
\begin{enumerate}
  \item[a)] les trois secteurs $v_1, v_2, v_3$ en $x$ sont dans trois
    composantes distinctes de $E \smallsetminus \Gamma$ ;
  \item[a')] $E \smallsetminus \Gamma$ a trois composantes ;
  \item[b)] $\Gamma_1, \Gamma_2, \Gamma_3$ sont des J-courbes ;
  \item[c)] $\Gamma_1$ et $\Gamma_2$ sont des J-courbes.
\end{enumerate}

\textbf{a) $\Leftrightarrow$ a').} Le côté d'un arc se transporte le long de
l'arc : les deux côtés de $\Gamma''_k$ sont, près de $x$, les secteurs $v_i$ et
$v_j$, et restent dans la même composante tout le long. Toute composante de $E
\smallsetminus \Gamma$ a un point adhérent à $\Gamma$, donc contient un côté
d'un arc, donc l'un des trois secteurs. Il y a ainsi au plus trois composantes,
et exactement trois si et seulement si les secteurs sont dans des composantes
deux à deux distinctes.

\textbf{b) $\Rightarrow$ c)} est trivial.

\textbf{c) $\Rightarrow$ a).} Deux points proches de $x$, de part et d'autre
de $\Gamma'_3$, sont dans $v_1$ et $v_2$. Comme $\Gamma'_3 \subset \Gamma_1$ et
que $\Gamma_1$ est une J-courbe, ils sont dans des composantes distinctes de $E
\smallsetminus \Gamma_1$, et a fortiori de $E \smallsetminus \Gamma$, dont les
composantes sont plus fines. De même $\Gamma'_1 \subset \Gamma_2$ sépare $v_2$
de $v_3$, et $\Gamma'_2 \subset \Gamma_1$ sépare $v_1$ de $v_3$\footnote{C'est
l'argument de la page, dont plusieurs mots sont illisibles ; l'indice de
« $E - \Gamma_2$ » dans la première phrase est d'une lecture incertaine. La
page dit « régions » pour composantes connexes.}.

\textbf{a) $\Rightarrow$ b).} La page s'arrête sur « Prouvons que $\Gamma_3$
(p.\ ex.) est courbe de Jordan » ; les formules de la page 2 contiennent la
suite\footnote{La rédaction de cette implication est de cette lecture ; elle
n'utilise que les formules des pages 2 et 16.}. Sous a), soient $V_1, V_2, V_3$
les trois composantes, $V_k \supset v_k$. Chaque $\Gamma''_k$ a $V_i$ d'un côté
et $V_j$ de l'autre, donc $U_k = V_i \cup \Gamma''_k \cup V_j$ est un ouvert
connexe, et
\[
E \smallsetminus \Gamma_k = U_k \sqcup V_k
\]
est la réunion disjointe de deux ouverts connexes non vides : ce sont ses deux
composantes. En tout point de $\Gamma''_i$ ($i \neq k$), les deux côtés sont
dans $V_k$ et dans $V_j \subset U_k$ ; en $x$ et $y$ de même. Donc $\Gamma_k$
est une J-courbe, et l'on a les relations de la page 16\footnote{La page 16 écrit $V_1 = U_1 \cap U_2$, comme pour $V_3$ ; c'est $U_2
\cap U_3$.} :
\[
V_k = U_i \cap U_j, \qquad U_k = V_i \cup V_j \cup \Gamma''_k
\qquad (\{i, j, k\} = \{1, 2, 3\}).
\]


Lue dans le plan, l'équivalence est une forme du \emph{théorème du thêta} : le
complémentaire d'un thêta plan a trois composantes, de frontières les trois
courbes $\Gamma_k$\footnote{Le nom est d'usage dans les exposés du théorème de
Jordan et dans la littérature sur les graphes plongés (« theta-curve
theorem ») ; les pages ne l'emploient pas.}. Mais elle est écrite pour servir de
\emph{pas} : elle transmet la propriété de deux courbes à la troisième. Sur la
façon dont ce pas démontre le théorème de Jordan pour les polygones, voir la
section sur le fil du dossier et la page 18.

\pagerange{16}{16}
\subsection*{Un ensemble d'indices, et une courbe dans une surface (page 16)}

Le haut de la page reprend le thêta avec des arcs $\Gamma'_\alpha$ indexés par
un ensemble $I$, tous joignant les deux points de $\partial$, un ensemble
$B_\alpha$ à deux éléments attaché à chaque arc, $B = \coprod_i B_i \simeq I
\times \{0, 1\}$, et, pour chaque $i$, $U_i$ la région de $\Gamma_i$ qui
contient $\Gamma''_i$, avec\footnote{La page écrit $U_j$ sous la réunion et
$V_j$ sous l'intersection ; on lit $V_j$ et $U_j$, ce qui redonne pour $I =
\{1, 2, 3\}$ les formules qui suivent sur la même page. Les mots qui disent ce
qu'est $B_\alpha$ sont en partie illisibles (« l'une des 2 origines, des pts
$x$ ») ; « ensemble » et « déterminent » sont d'une lecture incertaine.}
\[
U_i = \Gamma''_i \cup \bigcup_{j \neq i} V_j, \qquad
V_i = \bigcap_{j \neq i} U_j .
\]

Pour plus de trois arcs, $\Gamma_i$ n'est plus une courbe fermée, et la page ne
dit pas ce qu'elle désigne ; elle revient aussitôt au cas $I = \{1, 2, 3\}$.
Les lettres $B_i$, $U_i$, $V_i$ sont celles du cadre général des pages 7 et 5,
où $\dot{K} = \coprod_i B_i$ ; c'est sans doute pour ce cas que le cadre est
fait, mais la page ne le dit pas.

Le bas de la page énonce, pour une courbe fermée simple $C$ dans une surface
connexe $E$, l'équivalence de :
\begin{enumerate}
  \item[a)] $\operatorname{card} \pi_0(E \smallsetminus C) = 2$ ;
  \item[b)] il existe un ouvert $U \subset E$ tel que $\dot{U} = C$ et $V = E
    \smallsetminus \overline{U} \neq \emptyset$, c'est-à-dire $U \neq E
    \smallsetminus C$ ;
  \item[c)] pour un petit disque $D_x$ centré en un point $x \in C$, et deux
    points $x', x''$ de $D_x \smallsetminus C$ de part et d'autre de $C$,
    $x'$ et $x''$ ne sont pas dans la même composante de $E \smallsetminus
    C$\footnote{La page écrit « $U \subset C$ », « $V = U - \overline{U}$ »,
    « $\overline{U} \neq E \smallsetminus C$ » et, dans c), « $x'$ et $y$ » ;
    on lit $U \subset E$, $V = E \smallsetminus \overline{U}$, $U \neq E
    \smallsetminus C$ et $x''$. La parenthèse « $\{x, y\} \to \pi_0(X - C)$
    est surjectif » dit c) sous la forme : les deux côtés locaux atteignent
    toutes les composantes. La page écrit aussi $X$ pour $E$.}.
\end{enumerate}
La page note qu'on a toujours $\operatorname{card} \pi_0(E \smallsetminus C)
\leq 2$, et en marge qu'alors $U$ est une composante de $E \smallsetminus C$.
Tout cela est juste, par l'argument de la page 13 : $C$ a localement deux
côtés\footnote{Pour une courbe polygonale c'est immédiat ; pour une courbe
fermée simple quelconque dans une surface, c'est le théorème de Schoenflies,
local.} ; un côté se transporte le long de $C$ ; toute composante contient un
côté. D'où au plus deux composantes, et deux exactement si et seulement si les
deux côtés en un point sont séparés, ce qui est c). Sous b), $U$ et $V$ sont
deux ouverts disjoints non vides recouvrant $E \smallsetminus C$, d'où a) ;
sous a), l'une des composantes convient pour $U$. Sur une bande de Möbius,
l'âme est une courbe à un seul côté global, et $E \smallsetminus C$ est
connexe : la majoration est atteinte ou non selon la courbe\footnote{L'exemple
est de cette lecture.}.

\pagerange{5}{11}
\section*{II. Ouverts complémentaires, pseudo-partitions, recollement (pages 9, 11, 7, 5)}

Ces quatre feuillets portent une suite numérotée de sa main, 1) et 3) page 9,
la suite de 3) et 4) page 11, 5) page 7, la fin de 5) et 6) page 5 ; le 2)
n'est pas dans le dossier. On suit ses numéros. Dans toute la section, $X$ est
un espace topologique.

\pagerange{9}{9}
\subsection*{1) Deux ouverts dont chacun est l'extérieur de l'autre (page 9)}

Pour deux ouverts $U, V$ de $X$, les conditions suivantes sont équivalentes :
\begin{enumerate}
  \item[a)] $U = \complement \overline{V}$ et $V = \complement \overline{U}$ ;
  \item[b)] $U \cap V = \emptyset$ et $\complement(U \cup V) \subset
    \overline{U} \cap \overline{V}$, avec alors égalité ;
  \item[c)] $U \cap V = \emptyset$, $\overline{U \cup V} = X$ et $\dot{U} =
    \dot{V}$ ;
  \item[d)] $\overline{U} \cap \overline{V} = \dot{U} = \dot{V}$ et $U \cup V$
    est dense dans $X$\footnote{La page ne demande pas la densité dans d), et
    sans elle d) n'équivaut pas aux autres : dans la somme disjointe de
    $\mathbb{R}$ et d'un point isolé $p$, les ouverts $U = \left]-\infty,
    0\right[$ et $V = \left]0, +\infty\right[$ vérifient $\overline{U} \cap
    \overline{V} = \dot{U} = \dot{V} = \{0\}$, mais $\complement\overline{V}$
    contient $p$. L'exemple est de cette lecture.}.
\end{enumerate}
Deux ouverts disjoints sont disjoints de leurs adhérences mutuelles, ce qui
donne l'égalité de b) et, de proche en proche, les équivalences. Sous ces
conditions $U$ et $V$ se déterminent l'un l'autre ; et, $U$ étant donné, un
tel $V$ existe si et seulement si $\complement \overline{U}$ est dense dans
$\complement U$, c'est-à-dire si et seulement si $U = \operatorname{int}
\overline{U}$. Ce sont les \emph{ouverts réguliers}, et $V$ est le
pseudo-complément de $U$ dans l'algèbre de Boole qu'ils forment\footnote{Les
noms (ouvert régulier, algèbre de Boole des ouverts réguliers) sont classiques
bien avant 1976 ; la page n'en emploie aucun.}. Le modèle est celui d'une
courbe de Jordan : $U$ et $V$ sont le dedans et le dehors, et $\dot{U} =
\dot{V}$ est la courbe.

\subsection*{3) Un fermé partagé en deux pièces (page 9)}

Soient $K$ un fermé de $X$ et $U_K, U'_K$ deux ouverts relatifs de $K$ formant
dans $K$ une paire comme en 1) :
\[
\overline{U_K} = K \smallsetminus U'_K, \qquad \overline{U'_K} = K
\smallsetminus U_K, \qquad C = \overline{U_K} \cap \overline{U'_K}
\]
(les adhérences dans $K$ et dans $X$ coïncident, $K$ étant fermé), de sorte
que $K = U_K \sqcup C \sqcup U'_K$ et que $C$, fermé, est la frontière commune
des deux pièces dans $K$. Soit $B = U_K \cap \dot{K}$ la partie de la pièce
$U_K$ qui est sur le bord de $K$ ; elle est ouverte dans $\dot{K}$ et disjointe
de $C$. Posons
\[
U = U_K \smallsetminus B = \operatorname{Int}(U_K), \qquad
V = X \smallsetminus \overline{U_K} = \operatorname{Int}(X_1), \quad X_1 = X
\smallsetminus U_K .
\]
On a toujours $U = \complement \overline{V}$ : un point de $C$ ou de $B$ n'est
pas intérieur à $\overline{U_K}$, puisqu'il adhère à $U'_K$ ou à $\complement
K$. Pour avoir la paire complète, il faut de plus que \emph{tout point de $B$
soit adhérent à $U$} — ce qui a lieu si $K$ est l'adhérence de son
intérieur\footnote{L'hypothèse n'est pas sur la page, et sans elle la
conclusion est fausse. Dans le plan, soit $K$ le disque unité fermé augmenté du
segment $[1, 2] \times \{0\}$, $U_K$ la partie de $K$ où $y > 0$ ou $x > 1$,
$U'_K$ celle où $y < 0$ ; alors $C = [-1, 1] \times \{0\}$, $B$ contient le
segment $\left]1, 2\right] \times \{0\}$, qui est dans $\complement
\overline{U}$ mais pas dans $V$. L'exemple est de cette lecture.}. Alors
\[
V = \complement \overline{U}, \qquad
\dot{U} = \dot{V} = \overline{U} \cap \overline{V} = \complement(U \cup V) = B
\sqcup C, \qquad V \subset X_1 \subset \overline{V} ,
\]
et $(U, V)$ est une paire au sens de 1).

\emph{Corollaire.} $X_1 = \complement U_K$ est localement fermé si et
seulement si $B$ est fermé dans $X$, c'est-à-dire $\overline{B} \cap C =
\emptyset$ ; il suffit pour cela que $\dot{K} \cap C = \emptyset$, autrement dit
que $\dot{K} \subset U_K \cup U'_K$. En effet $\overline{X_1} \smallsetminus X_1
= B$, et une partie est localement fermée si et seulement si son adhérence
moins elle-même est fermée ; et $\overline{B} \subset \dot{K} \cap (B \cup
C)$\footnote{La page dit d'abord $B$ « relativement fermé dans $X$ », lecture
incertaine, puis, après plusieurs ratures, que $B$ doit être fermé dans $B \cup
C$, et suppose $\overline{B} \cap C = \emptyset$ ; la fin de la page, très
raturée, n'est lue qu'en partie. Le « $V_K$ » de sa dernière ligne est
$U'_K$.}.

\pagerange{11}{11}
\subsection*{3, suite) et 4) Trois ouverts et trois murs (page 11)}

Supposons $\dot{K} \cap C = \emptyset$, faisons de même pour l'autre pièce,
$B' = U'_K \cap \dot{K}$, $U' = \operatorname{Int}(U'_K)$, $V' = X
\smallsetminus \overline{U'_K}$, et supposons $B' \subset \overline{U'}$ comme
$B \subset \overline{U}$. Soit $W = V \cap V' = \complement K$. Alors $X$ est
réunion disjointe
\[
X = U \sqcup U' \sqcup W \sqcup B \sqcup B' \sqcup C ,
\]
de trois ouverts et de trois fermés, avec\footnote{La première égalité de la
deuxième ligne commence sur la page par $\overline{U^{*}} \cap
\overline{W}$, lapsus probable.}
\[
\left\{
\begin{array}{l}
\overline{U} \cap \overline{U'} = C, \quad \overline{U} \cap \overline{W} = B,
\quad \overline{W} \cap \overline{U'} = B', \\
\dot{U} = B \cup C, \quad \dot{U}' = B' \cup C, \quad \dot{W} = B \cup B', \\
K = \overline{U \cup U'} = \complement W, \quad \operatorname{int} K = U \cup
U' \cup C.
\end{array}
\right.
\]
La figure de la page
est celle d'un disque $K$ coupé en deux par un mur intérieur $C$ qui ne touche
pas le bord, dans une couronne : on peut penser à $X$ privé de quelques points,
le mur $C$ allant d'un trou à l'autre.

La page appelle cela une \emph{pseudo-partition} sans définir le mot. On le
définit ainsi\footnote{La définition est de cette lecture, tirée des relations
que la page vient d'écrire ; c'est sous elle que l'équivalence 4) est
vérifiée.} : six parties $U, U', W, B, B', C$ forment une pseudo-partition de
$X$ si elles le partagent, si $U, U', W$ sont ouverts, et si
\[
\overline{U} = U \sqcup B \sqcup C, \qquad \overline{U'} = U' \sqcup B' \sqcup
C, \qquad \overline{W} = W \sqcup B \sqcup B' .
\]

\textbf{4).} Pour trois ouverts $U, U', W$ de $X$, sont équivalentes :
\begin{enumerate}
  \item[a)] ils sont deux à deux disjoints, de réunion dense,
    $\dot{U} \subset \overline{U' \cup W}$, $\dot{U}' \subset \overline{U \cup
    W}$, $\dot{W} \subset \overline{U \cup U'}$, et $\overline{U} \cap
    \overline{U'} \cap \overline{W} = \emptyset$ ;
  \item[b)] avec $C = \overline{U} \cap \overline{U'}$, $B = \overline{U} \cap
    \overline{W}$, $B' = \overline{W} \cap \overline{U'}$, les six parties
    forment une pseudo-partition de $X$.
\end{enumerate}
Sous a), les trois murs sont deux à deux disjoints (leurs intersections sont
dans le point triple, vide), disjoints des trois ouverts, et ils recouvrent
le reste : un point de $\overline{U} \smallsetminus U$ est dans
$\overline{U'}$ ou dans $\overline{W}$. Réciproquement b) donne a) en lisant
les adhérences. C'est le dessin de trois pays deux à deux voisins dont les
trois frontières ne se rencontrent pas.

Le thêta de la page 2, privé de ses deux sommets, en est l'exemple : dans $X = E
\smallsetminus \partial$, les trois régions $V_1, V_2, V_3$ et les trois arcs
ouverts $\Gamma''_1, \Gamma''_2, \Gamma''_3$ forment une pseudo-partition, le
point triple $\partial$ ayant été retiré\footnote{Le rapprochement est de
cette lecture ; il suppose a) de la page 13.}.

La page revient enfin (« Je reviens aux données ») sur le lien entre les deux
situations. Elles se déterminent l'une l'autre : d'une pseudo-partition on tire
$K = \complement W$, $U_K = U \cup B$, $U'_K = U' \cup B'$, qui vérifient les
hypothèses de 3), avec $C \cap \dot{K} = \emptyset$ puisque $\dot{K} = \dot{W}
= B \cup B'$, et la page le note sous la forme $C = K \smallsetminus (U_K \cup
U'_K)$\footnote{La première phrase de ce paragraphe est en partie illisible ;
la réciprocité est rédigée ici par cette lecture.}.

\pagerange{7}{7}
\subsection*{5) Découper le long des murs intérieurs (page 7)}

On généralise à un nombre quelconque de pièces. Soient $K$ un fermé de $X$ et
$(U_i)_{i \in I}$ une famille d'ouverts relatifs de $K$, deux à deux disjoints,
tels que $\dot{K} \subset \bigcup_i U_i$. Posons $B_i = \dot{K} \cap U_i$, de
sorte que $\dot{K} = \coprod_i B_i$ est une somme topologique, chaque $B_i$
étant ouvert, donc aussi fermé, dans $\dot{K}$. Soit $C_i = \overline{U_i}
\smallsetminus U_i$ la frontière relative de $U_i$ dans $K$ ; on suppose
$C_i \subset \overline{K \smallsetminus \overline{U_i}}$ (le mur $C_i$ est
adhérent à ce qui est de l'autre côté), et, comme en 3), que $B_i$ est adhérent
à $U_i \smallsetminus B_i$\footnote{Cette seconde hypothèse n'est pas sur la
page ; c'est celle qui manquait déjà en 3).}. Enfin $W = X \smallsetminus K$.

\textbf{Le complémentaire d'une pièce.} $X \smallsetminus U_i$ est localement
fermé, et
\[
\operatorname{Int}(X \smallsetminus U_i) = X \smallsetminus \overline{U_i} = (X
\smallsetminus U_i) \smallsetminus C_i, \qquad \overline{X \smallsetminus U_i}
= (X \smallsetminus U_i) \sqcup B_i ;
\]
les deux ouverts $\operatorname{Int}(X \smallsetminus U_i)$ et
$\operatorname{Int}(U_i) = U_i \smallsetminus B_i$ forment une paire au sens de
1)\footnote{La page écrit $\overline{X_i} = X_i \cup B'_i$ avec $B'_i = \dot{K}
\smallsetminus B_i$ : mais $B'_i$ est déjà dans $X \smallsetminus U_i$, et ce
que l'adhérence ajoute est $B_i$. Elle écrit aussi « $U_i \smallsetminus C_i$ »,
qui est $U_i$ et n'est pas ouvert dans $X$ ; on lit $U_i \smallsetminus B_i$.
L'indice de $U_i$ est surchargé à chaque occurrence de la page, un accent
($U'_i$) étant lisible par endroits.}.

\textbf{La pièce recollée.} Posons maintenant $X_i = W \cup \overline{U_i}$ :
l'extérieur de $K$ et la seule pièce $i$, avec son mur. L'inclusion $W \subset
X_i$ est une immersion ouverte, la frontière de $W$ dans $X_i$ est $B_i$, et
$X_i$ est la somme amalgamée de deux fermés le long de leur intersection\footnote{La page désigne par la
même lettre $X_i$ le complémentaire $X \smallsetminus U_i$ et cette pièce
recollée ; les deux ne coïncident pas en général. Quand il n'y a que deux
pièces, $K = U_1 \sqcup C \sqcup U_2$ comme en 3), on a $X \smallsetminus U_1
= W \cup \overline{U_2}$ : la double notation est un échange des deux indices.
L'égalité $\overline{W} \cap X_i = W \cup B_i$ demande que $C_i \cap \dot{K} =
\emptyset$, voir plus bas. La dernière ligne de la page n'est lue qu'en
partie.} :
\[
X_i = (W \cup B_i) \amalg_{B_i} \overline{U_i} .
\]


\pagerange{5}{5}
\subsection*{5, suite) L'espace découpé, et 6) un lemme annoncé (page 5)}

On recolle alors toutes les pièces le long de l'extérieur, sans les recoller
entre elles : soit
\[
\mathcal{X} = \overline{W} \amalg_{\coprod_i B_i} \Bigl(\coprod_i
\overline{U_i}\Bigr) \simeq \coprod_{i\,(W)} X_i ,
\]
somme amalgamée des $X_i$ au-dessus de $W$\footnote{« peut aussi s'identifier
à » et « somme amalgamée » sont des lectures incertaines ; la formule, elle,
est nette.}. C'est $X$ découpé le long des murs intérieurs $C_i$ : deux pièces
qui se touchent dans $X$ le long d'un mur sont séparées dans $\mathcal{X}$. La
flèche évidente $\pi : \mathcal{X} \to X$ est l'identité sur chaque morceau.
L'image $\mathcal{X}_i$ de $X_i$ est ouverte dans $\mathcal{X}$, et $X_i \to
\mathcal{X}_i$ est un homéomorphisme. On en tire :
\begin{itemize}
  \item $\pi$ est un homéomorphisme local en tout point de $\mathcal{X}
    \smallsetminus \coprod_i C_i$ : un point de $B_i$ a dans $X$ un voisinage
    contenu dans $W \cup U_i$ ;
  \item en un point de $C_i$, $\pi$ est localement un plongement fermé, $X_i$
    coïncidant près de ce point avec $\overline{U_i}$\footnote{La page dit :
    « $\mathcal{X} \to X$ est étale en les pts de $\mathcal{X} - \coprod_i C_i$,
    i.e.\ une [illisible] \emph{fermée} en les pts de $\coprod C_i$ », ce
    dernier mot étant d'une lecture incertaine. Le « i.e. » ne relie pas deux
    énoncés équivalents ; on rédige les deux. Les vérifications sont de cette
    lecture.}.
\end{itemize}

La page se clôt sur une inquiétude : « vérifier si dans les développements
précédents, il n'y avait pas lieu d'exiger $\bigcup C_i \subset
\operatorname{int} K$, i.e.\ $C_i \cap \dot{K} = \emptyset$ ». L'hypothèse est
bien nécessaire — c'est elle qui fait de $B_i$ la frontière de $W$ dans $X_i$ —
mais elle est déjà contenue dans $\dot{K} \subset \bigcup_i U_i$ : un point de
$\dot{K}$ est dans un $U_j$, et $U_j$, ouvert relatif disjoint de $U_i$, est
disjoint de $\overline{U_i}$ et donc de $C_i$\footnote{La réponse est de cette
lecture. La marge porte « vérifier », suivi de deux mots illisibles.}.

\textbf{6) Lemme d'immersion de $\mathcal{X}$.} La page annonce un lemme sous
ce titre, puis remplace $K$, près de son bord, par un fermé plus petit $K^{*}
\subset K$, avec $U^{*}_i = U_i \cap K^{*}$, $U'^{*}_i = U'_i \cap K^{*}$, en
demandant que $\dot{K}^{*}$ soit dans les $U_i$ et que $C_i \subset
\overline{U^{*}_i} \cap \overline{U'^{*}_i}$, ce qui a lieu si
$\operatorname{int} K^{*} \supset \bigcup C_i$. Plusieurs conditions sont
barrées, le lemme n'est pas énoncé\footnote{« Lemme », « remplaçant » et « on
suppose » sont des lectures incertaines ; un mot illisible suit « de
$\mathcal{X}$ ». La parenthèse ouverte avant « $\operatorname{int} K'$ » n'est
pas fermée, et le $K'$ y est probablement le $K^{*}$ des lignes
précédentes.}.

Tout ce cadre est celui du « découpage », qu'on retrouve dans le dossier 81 :
la Version I de ce dossier (pages 13 à 24) découpe un espace connexe par des
coupures fermées à deux côtés et compte les morceaux, et elle observe que
l'exemple du disque coupé par des arcs est une forme du théorème de Jordan.

\pagerange{6}{12}
\section*{III. Un calcul abandonné : filtrations et découpage d'un complexe (pages 6 et 12)}

\pagerange{6}{6}
\subsection*{Une filtration indexée par un ensemble ordonné (page 6)}

Les deux pages sont barrées. Soit $E$ un ensemble ordonné, $E_0 = M(E)$
l'ensemble de ses éléments minimaux (les sommets), et pour $e \in E$, $E_e =
\{x \in E \mid x \leq e\}$, $M(e) = M(E_e)$ les sommets de $e$ ; $X = |E|$ est
la réalisation géométrique, réunion des cellules fermées $|E_e|$. On se donne
une application $\varphi : E_0 \to I$ dans un ensemble ordonné, et l'on filtre
$X$ par
\[
X_i = \bigcup_{\max \varphi(M(e)) \leq i} |E_e|, \qquad \partial X_i =
\bigcup_{j < i} X_j\footnote{La page 6 écrit « $\varphi(M(e)) \leq i$ » ; la
page 12 précise $\mu(\varphi(M(e)))$, qu'on lit comme le plus grand
élément.}.
\]
Pour un drapeau $d = (i_0 > \cdots > i_k)$ de $I$, la page considère des
espaces $\Sigma_d(\mathcal{X})$, avec $\Sigma_{\{i\}} = \Sigma(X_i, \partial
X_i)$, des flèches $\Sigma_d \to \Sigma_{d'}$ pour $d' \subset d$, et
l'ensemble $D(d) = \pi_0 \Sigma_d$ de leurs composantes connexes.

\pagerange{12}{12}
\subsection*{Les composantes, et le découpage d'un complexe (page 12)}

La page 12 énonce, pour $d' \subset d$, un carré commutatif de bijections
\[
\alpha_d : D(d) = \pi_0 \Sigma_d(\mathcal{X}) \xrightarrow{\ \sim\ } E(d) = \{e
\in E \mid \varphi(M(e)) = d\},
\]
compatibles aux restrictions $D(d) \to D(d')$ et $E(d) \to E(d')$. Ce que sont
$\Sigma$ et $\Sigma_d$, aucune des deux pages ne le dit ; la bijection ne peut
donc pas être vérifiée ici\footnote{La flèche verticale de gauche porte une
étiquette surchargée, lue $\mathrm{P}_{d',d}$ ; le sens d'une flèche du
diagramme de la page 6 est incertain.}.

Le bas de la page donne en revanche une définition, qui éclaire sans doute ce
que $\Sigma$ veut dire. Pour un complexe simplicial $K$, vu comme un ensemble
de sommets $K_0$ et un ensemble de simplexes, et un sous-complexe $L$, on pose
\[
\tilde{K} = \Sigma(K, L), \qquad \tilde{K}_0 = \varinjlim_{\sigma \in K
\smallsetminus L} \sigma \xrightarrow{\ \pi\ } K_0 ,
\]
la limite inductive des ensembles de sommets des simplexes qui ne sont pas dans
$L$, indexée par l'ensemble ordonné $K \smallsetminus L$ (stable par passage à
un simplexe plus grand). Deux simplexes hors de $L$ ne sont recollés que le
long de faces qui sont elles aussi hors de $L$ : c'est $K$ \emph{découpé le
long de} $L$. Par exemple, si $K$ est formé de deux arêtes $ab$ et $bc$ et $L =
\{b\}$, le sommet $b$ se dédouble et $\tilde{K}$ est formé de deux arêtes
disjointes\footnote{L'exemple est de cette lecture.}. La page se demande, sur
une ligne barrée et avec un « ? » au-dessus du signe $\sim$, si la réalisation
de $\tilde{K}$ est le découpage topologique $\Sigma(|K|, |L|)$, présenté comme
une limite inductive de simplexes. C'est l'analogue combinatoire de l'espace
$\mathcal{X}$ des pages 7 et 5.

\pagerange{10}{10}
\section*{IV. Une page d'algèbre (page 10)}

Feuillet plié, sans rapport avec le reste du dossier, écrit tête-bêche. Dans un
groupe engendré par $\lambda_0, \lambda_1, \lambda_2$, on considère
l'automorphisme
\[
t(\lambda_0) = \lambda_0, \qquad t(\lambda_1) = \lambda_1 \lambda_2
\lambda_1^{-1}, \qquad t(\lambda_2) = \lambda_1 .
\]
La page écrit aussi $t(\lambda_1) = \lambda_0^{-1}\lambda_2\lambda_0$, ce qui
vaut si $\lambda_0 \lambda_1 \lambda_2 = 1$, relation qu'on suppose
désormais\footnote{La relation n'est pas écrite ; on la déduit de cette
égalité et de « $\operatorname{int}(\lambda_1\lambda_2) =
\operatorname{int}(\lambda_0^{-1})$ ». C'est la présentation du groupe
fondamental du plan privé de deux points, libre sur $\lambda_1, \lambda_2$.}.
Alors $t$ préserve le produit $\lambda_1 \lambda_2$, et
\[
t^2(\lambda_1) = \operatorname{int}(\lambda_1 \lambda_2)(\lambda_1), \qquad
t^2(\lambda_2) = \lambda_1 \lambda_2 \lambda_1^{-1} =
\operatorname{int}(\lambda_1 \lambda_2)(\lambda_2),
\]
d'où $t^2 = \operatorname{int}(\lambda_1\lambda_2) =
\operatorname{int}(\lambda_0^{-1})$, où $\operatorname{int}(g)(h) = g h
g^{-1}$\footnote{La page intitule ce calcul $t^3$ et écrit « $t^3 =
\operatorname{int}(\lambda_0)$ » ; mais la ligne « $t^3(\lambda_2) =
t(\lambda_1)$ » montre que c'est $t \circ t$ qui est calculé, et le résultat
de la page est $\operatorname{int}(\lambda_0^{-1})$. Le premier facteur de la
deuxième ligne de la page, lu $t(q_\nu)$, est incertain.}. C'est l'action d'un
générateur du groupe de tresses à deux brins (l'action de Hurwitz), dont le
carré, le « tour complet », agit par conjugaison\footnote{Les noms sont de
cette lecture.}. À droite, un croquis au crayon pâle évoque l'induction
$G \times_H E$ d'un $H$-espace et une flèche $B_H \to B_G$ d'espaces
classifiants, avec $f^{*}$ et $f_{!}$ ; la lecture en est incertaine dans le
détail.

\pagerange{18}{18}
\section*{V. Recoller deux J-polygones (page 18)}

\textbf{Proposition.} Soient $\Gamma, \Gamma'$ deux J-polygones distincts du
plan $E$ tels que $\Gamma \cap \Gamma' = I$ soit un arc polygonal non réduit à
un point\footnote{« Distincts » n'est pas sur la page ; si $\Gamma = \Gamma'$,
l'arc $I$ n'en est pas un et $\Gamma''$ ci-dessous est vide. Deux polygones
distincts ne sont jamais contenus l'un dans l'autre (lemme de la page 31). Le
mot « du » devant « plan » est d'une lecture incertaine.}. On est dans l'un
des deux cas suivants, qui s'excluent :
\begin{enumerate}
  \item[a)] $\Gamma \smallsetminus I \subset \operatorname{Ext}\Gamma'$ et
    $\Gamma' \smallsetminus I \subset \operatorname{Ext}\Gamma$ ; de façon
    équivalente $\operatorname{Int}\Gamma' \subset \operatorname{Ext}\Gamma$ et
    $\operatorname{Int}\Gamma \subset \operatorname{Ext}\Gamma'$ ;
  \item[b)] $\Gamma \smallsetminus I \subset \operatorname{Ext}\Gamma'$ et
    $\Gamma' \smallsetminus I \subset \operatorname{Int}\Gamma$, ou l'inverse ;
    de façon équivalente $\operatorname{Int}\Gamma' \subset
    \operatorname{Int}\Gamma$ et $\operatorname{Ext}\Gamma \subset
    \operatorname{Ext}\Gamma'$, ou l'inverse.
\end{enumerate}
Soit $\Gamma'' = (\Gamma \cup \Gamma') \smallsetminus (I \smallsetminus
\partial I) = \overline{\Gamma \smallsetminus I} \cup \overline{\Gamma'
\smallsetminus I}$, réunion de deux arcs qui ne se rencontrent qu'en leurs
extrémités $\partial I$ : c'est un polygone. C'est un J-polygone, et :
\begin{itemize}
  \item dans le cas a), $\operatorname{Int}\Gamma'' = \operatorname{Int}\Gamma
    \cup \operatorname{Int}\Gamma' \cup (I \smallsetminus \partial I)$ et
    $\operatorname{Ext}\Gamma'' = \operatorname{Ext}\Gamma \cap
    \operatorname{Ext}\Gamma'$ ;
  \item dans le cas b), $\operatorname{Int}\Gamma'' = \operatorname{Int}\Gamma
    \cap \operatorname{Ext}\Gamma'$ et $\operatorname{Ext}\Gamma'' =
    \operatorname{Ext}\Gamma \cup \operatorname{Int}\Gamma' \cup (I
    \smallsetminus \partial I)$.
\end{itemize}
Chaque fois, la page écrit la partition de $E$ en sept parties qui le
justifie : dans le cas a), $\operatorname{Int}\Gamma$,
$\operatorname{Int}\Gamma'$, $\Gamma \smallsetminus I$, $\Gamma'
\smallsetminus I$, $I \smallsetminus \partial I$, $\partial I$ et
$\operatorname{Ext}\Gamma \cap \operatorname{Ext}\Gamma'$ ; dans le cas b), la
même avec $\operatorname{Ext}\Gamma$ et $\operatorname{Int}\Gamma \cap
\operatorname{Ext}\Gamma'$ à la place des deux parties qui changent. Les deux
ouverts proposés pour $\operatorname{Int}\Gamma''$ et
$\operatorname{Ext}\Gamma''$ sont connexes, la partie $I \smallsetminus
\partial I$ soudant les deux régions qu'elle sépare.

La dichotomie, que la page énonce sans la démontrer, vient de ce que $\Gamma
\smallsetminus I$ est connexe et disjoint de $\Gamma'$, donc dans
$\operatorname{Int}\Gamma'$ ou dans $\operatorname{Ext}\Gamma'$, et de même pour
$\Gamma' \smallsetminus I$ ; le quatrième cas est impossible, car si $\Gamma'
\smallsetminus I \subset \operatorname{Int}\Gamma$, alors
$\operatorname{Ext}\Gamma$, connexe, non borné et disjoint de $\Gamma'$, est
dans $\operatorname{Ext}\Gamma'$, et $\Gamma \smallsetminus I$, adhérent à
$\operatorname{Ext}\Gamma$, y est aussi\footnote{L'argument est de cette
lecture.}.

C'est l'énoncé qu'il faut pour achever la récurrence décrite dans le fil du
dossier : si l'on coupe un polygone $\Gamma$ par une diagonale $[s, s']$ en deux
polygones $\Lambda^1, \Lambda^2$ qui se rencontrent le long de $I = [s, s']$,
alors $\Gamma = \Lambda''$, et la proposition dit qui sont son intérieur et son
extérieur. Le cas b) est celui d'une diagonale extérieure — l'un des deux
polygones contient l'autre —, le cas a) celui d'une diagonale intérieure. Ce
sont exactement les deux situations des §§ 3 et 4 du mémoire (lemmes XIV et XV,
lemmes XVIII et XIX, ci-dessous)\footnote{Le rapprochement est de cette
lecture.}.

\pagerange{20}{47}
\section*{VI. Le mémoire du 15 juillet 1976, et sa lecture (pages 20 à 47)}

Les pages 20 à 47 sont, pour l'essentiel, un mémoire d'une autre main, à
l'encre bleue, paginé de I à XIX, signé et daté « le 15 juillet 1976,
Montpellier » à sa dernière page. Il n'est pas recomposé dans cette édition, et
cette lecture n'en dit pas plus que les notes de la transcription : ses énoncés
sont nommés comme elles les nomment, ses démonstrations ne sont pas données,
son auteur n'est pas nommé\footnote{C'est la règle de la transcription pour le
texte d'autrui, reprise du dossier 100, et celle de RIGHTS.md.}. Ce qui est lu
ici, ce sont les marques de Grothendieck, ses feuillets intercalés, et la
mathématique des énoncés.

\pagerange{20}{22}
\subsection*{Les deux théorèmes, et les définitions (pages 20 à 22)}

Le titre est « Le théorème de Jordan et l'existence de trois bons sommets pour
un polygone de $\mathbb{R}^2$ ». \emph{Théorème 1} (Jordan) : tout polygone
est un J-polygone. \emph{Théorème 2}, attribué avec un point d'interrogation à
Dehn, « non publié » : tout polygone a au moins trois bons sommets, dont deux au
moins à triangle intérieur. Un polygone est la réalisation dans le plan d'un
polygone abstrait connexe, de sommets $S(\Gamma)$ et d'arêtes $A(\Gamma)$ — ce
que le dossier 89 (pages 19 et 20) appelle un polygone combinatoire, réalisé.
La \emph{définition 1} (page 22) dit ce qu'est un J-polygone : $\mathbb{R}^2
\smallsetminus \Gamma$ est réunion de deux ouverts connexes, l'un relativement
compact, l'autre non ; et un petit disque centré sur $\Gamma$ est coupé en deux
parties connexes, l'une à l'intérieur, l'autre à l'extérieur. La
\emph{définition 2} : un bon sommet est un sommet non aligné avec ses deux
voisins, dont le triangle ne contient aucun autre sommet de $\Gamma$.

Aujourd'hui, un sommet dont le triangle est intérieur au polygone et ne
contient pas d'autre sommet est une \emph{oreille}, et un sommet dont le
triangle est extérieur une \emph{bouche}. Le théorème des deux oreilles est de
G.~Meisters (1975) ; l'existence d'une bouche pour un polygone non convexe est
attribuée à G.~Toussaint (1991). Le théorème 2 du mémoire est la conjonction des
deux, un polygone convexe ayant au moins trois oreilles. Un manuscrit inédit de
Max Dehn sur le théorème de Jordan pour les polygones a été publié par
H.~Guggenheimer en 1977\footnote{G.~Meisters, « Polygons have ears »,
\emph{Amer. Math. Monthly} 82 (1975) ; G.~Toussaint, « Anthropomorphic
polygons », \emph{Amer. Math. Monthly} 98 (1991) ; H.~Guggenheimer, « The
Jordan curve theorem and an unpublished manuscript by Max Dehn », \emph{Arch.
Hist. Exact Sci.} 17 (1977). Références de mémoire, non vérifiées. Ni le mémoire
ni ses marges ne citent Meisters ; on n'a pas vérifié si le manuscrit de Dehn
contient l'énoncé que le mémoire lui attribue.}.

Ses marques sur ces pages sont des exigences de précision. Dans la définition
1, il ajoute « non vides ! » aux deux parties du disque — sans quoi la
condition ne dirait rien — et note que l'extérieur est « non rel.\ compact ».
Il ajoute « et non vide » à la définition du polygone, « de $\mathbb{R}^2$ »
où l'ambiant manquait. À la remarque 2 (il suffit qu'un tel disque vérifie la
condition), il donne la raison : si $D_x \subset D'_x$, l'application $\pi_0(D_x
\smallsetminus D_x \cap \Gamma) \to \pi_0(D'_x \smallsetminus D'_x \cap
\Gamma)$ est bijective. Au pied de la page 20, il écrit « faux » en face de la
fin de la remarque 1, qui affirme que deux points $x, y$ d'un polygone le
coupent en « un seul couple d'arcs » ; la feuille suivante dit
pourquoi\footnote{La remarque 1 se poursuit page 22, que la transcription
résume sans en citer la fin ; que le « faux » vise l'unicité du couple est une
lecture, appuyée sur le NB de la page 21. Les marques en pointillé et le mot
« plan » au crayon, page 20, sont d'une main non identifiée ; le « or.\ ? »
à l'encre noire, en face de « 1-complexe connexe », peut demander si les arcs
sont orientés.}.

\textbf{Ses notes, page 21.} Pour $x \neq y$ sur $\Gamma$, les deux arcs
$\Gamma^1_{xy}, \Gamma^2_{xy}$ vérifient
\[
\Gamma^1_{xy} \cap \Gamma^2_{xy} = \{x, y\}, \qquad \Gamma^1_{xy} \cup
\Gamma^2_{xy} = \Gamma, \qquad S'(\Gamma^1_{xy}) \sqcup S'(\Gamma^2_{xy}) =
S(\Gamma) \smallsetminus \{x, y\},
\]
où $S'$ désigne les sommets intérieurs d'un arc\footnote{La troisième égalité
suppose $x, y$ sommets, ou rendus tels par l'opération a) ; quelques signes
serrés au bout de la ligne ne sont pas déchiffrés.}. \textbf{NB} : ce n'est pas
le couple $(\Gamma^1_{xy}, \Gamma^2_{xy})$ qui est déterminé, mais la paire
$\{\Gamma^1_{xy}, \Gamma^2_{xy}\}$. Suivent les deux opérations dont tout le
reste est fait : a) ajouter des sommets ; b) diviser un polygone par deux
sommets distincts en deux lignes polygonales, l'opération inverse étant de
recoller deux lignes polygonales de mêmes extrémités\footnote{« rajouter » et
« mêmes » sont des lectures incertaines.}.

\pagerange{23}{28}
\subsection*{Voisinages d'un arc, ensembles convexes (pages 23 à 28)}

Le \emph{lemme I} (page 24) donne, pour un arc d'un J-polygone, un système
fondamental de voisinages ouverts $V$ tels que $V \cap \operatorname{int}\Gamma$
et $V \cap \operatorname{ext}\Gamma$ soient connexes. Il précise l'énoncé
(« $x \neq y$ », « ouverts »), relève d'un « ! » l'appel à l'axiome du choix,
et d'un « ? » l'aveu « on utilise le fait, sans le dire, que deux points d'un
ouvert connexe du plan peuvent être reliés par un arc de polygone ».

Le § 2 (pages 26 et 28) rappelle les ensembles convexes : le \emph{lemme II}
sur le bord de l'enveloppe convexe d'une partie (il ajoute « finie »), une
\emph{convention} qui réserve le mot « ensemble convexe » à l'enveloppe convexe
d'un ensemble fini de points, les \emph{lemmes III, IV, V} sur les droites
d'appui et la visibilité, la \emph{définition 3} (un polygone convexe est le
bord d'un ensemble convexe) et le \emph{lemme VI} (un polygone convexe est un
J-polygone, et tout sommet non aligné avec ses voisins y est un bon sommet). Ses
marges : « jargon ? », « !!! » devant la convention, « C'est une salade… » sur
l'intervention du lemme II dans la compacité, « vis.\ ? », et, au lemme V, la
précision que $s, s'$ sont sur une arête de $C$ et que $x$ et $C$ sont de part
et d'autre de leur droite. Les pages 23 et 27 ne portent que des dessins de sa
main : un polygone étoilé doublé d'un contour parallèle, et un croquis qui
paraît illustrer la démonstration du lemme IV.

\pagerange{29}{36}
\subsection*{Diagonales extérieure et intérieure (pages 29 à 36)}

La démonstration du théorème 1 repose sur l'existence d'une diagonale. Il y a
deux façons d'en trouver une, et le mémoire et ses feuillets à lui les
traitent toutes deux.

\textbf{Une diagonale extérieure, sans hypothèse (page 31).} Le mémoire établit
(page 30) que l'enveloppe convexe $E(\Gamma)$ d'un polygone est celle de ses
sommets et que les sommets de son bord $\partial E(\Gamma)$ sont des sommets de
$\Gamma$ ; puis (page 32) le \emph{lemme VII}, $\operatorname{ext}\partial
E(\Gamma) \subset \operatorname{ext}\Gamma$ pour un J-polygone, et le
\emph{lemme VIII} : un J-polygone non convexe a deux sommets $s, s'$ avec $s'$
visible de $s$ par l'extérieur. Grothendieck met le « J- » entre parenthèses et
écrit « inutile » en face. Sa page 31 dit pourquoi.

\emph{Lemme.} Deux polygones distincts $\Gamma, \Gamma'$ ne sont pas contenus
l'un dans l'autre, et $\Gamma \cap \Gamma'$ est une réunion finie d'arcs
disjoints, éventuellement réduits à un point. Si $\Gamma$ n'est pas convexe et
$\Gamma' = \partial E(\Gamma)$, l'ensemble $\Phi$ des extrémités des
composantes de $\Gamma \cap \Gamma'$ est contenu dans $S(\Gamma)$, et pour tout
$s \in \Phi$ il existe $t \in \Phi$, $t \neq s$, tel que $]s, t[\, \subset
\Gamma' \smallsetminus \Gamma$ ; alors $]s, t[\, \subset
\operatorname{Ext}\Gamma$. Un tel $t$ est unique, sauf si $s$ est un point isolé
de $\Gamma \cap \Gamma'$, auquel cas il y en a exactement deux.

\emph{Corollaire.} Un polygone non convexe a deux sommets distincts $s, s'$
tels que $]s, s'[\, \subset \operatorname{Ext}\Gamma \cap \partial E(\Gamma)$.

\textbf{NB} : $\operatorname{Ext}\Gamma$, la composante non bornée de
$\mathbb{R}^2 \smallsetminus \Gamma$, est défini pour tout polygone, et pas
seulement pour un J-polygone.

Tout cela est juste, et se voit ainsi. Un point de $\Gamma \cap \Gamma'$ qui
n'est pas un sommet de $\Gamma$ est sur un segment de $\Gamma$ contenu dans
$E(\Gamma)$ et passant par un point de son bord : ce segment est dans une droite
d'appui, donc dans $\Gamma'$ de part et d'autre du point, qui n'est pas une
extrémité. Une composante de $\Gamma' \smallsetminus \Gamma$ est un arc ouvert
de $\Gamma'$ ; elle ne contient aucun sommet de $\Gamma'$, puisque ceux-ci sont
des sommets de $\Gamma$, et c'est donc un segment ouvert $]s, t[$, dont les
extrémités sont dans $\Phi$. Enfin $]s, t[$ est connexe, disjoint de $\Gamma$, et
adhérent au complémentaire de $E(\Gamma)$, qui est connexe, non borné et
disjoint de $\Gamma$ : il est dans $\operatorname{Ext}\Gamma$. Rien n'y fait
appel au théorème de Jordan\footnote{Les vérifications sont de cette lecture.
Sur la page, « non convexe » et « l'ens.\ $\Phi$ des » sont d'une lecture
incertaine, et « tt » et « néc. » dans le NB aussi. La formule de la
transcription « $s$ $]s, s'[$ » du corollaire porte un mot illisible barré.}.
C'est aussi ce que le mémoire, page 36, énonce comme « corollaire à la
démonstration du lemme VIII », pour tout polygone ; ses marges en face disent
« !! ».

Combinée avec le pas du thêta de la page 13, cette diagonale suffit au
théorème 1 : si $[s, s']$ est une diagonale extérieure, $\Gamma \cup [s, s']$
est un thêta dont les deux autres courbes fermées ont moins de sommets que
$\Gamma$\footnote{Le rapprochement est de cette lecture ; voir le fil du
dossier.}.

\textbf{Une diagonale intérieure (pages 29, 33, 35 et 36).} Le mémoire, page
33, établit qu'en un sommet d'un polygone convexe l'angle intérieur est au plus
$\pi$ (\emph{lemme IX}, « l'analogue de la positivité de la courbure »), et
qu'un J-polygone a au moins trois sommets d'angle intérieur $< \pi$
(\emph{lemme X}) — les sommets de $\partial E(\Gamma)$ en fournissent. Le
\emph{lemme XI} (page 35) affirme qu'un J-polygone a deux sommets $s, s'$ avec
$s'$ visible de $s$ par l'intérieur ; on prend un sommet $s$, ses voisins $s',
s''$ et le triangle $T$ de sommets $s, s', s''$.

Ses objections sont précises, et toutes fondées. Au crayon : « faux pour $n = 3$
! » — un triangle n'a pas de diagonale. À l'encre : il ajoute « à $n \geq 3$
côtés », « OPS $\Gamma$ non convexe (sinon c'est clair) ». En face de « soit
$s$ un sommet satisfaisant la propriété énoncée au lemme VIII », il écrit
« Lemme X », et au crayon : « il faut prendre $s$ comme dans la \emph{dém.} du
lemme VIII (pas son énoncé !) » — il faut que l'angle intérieur en $s$ soit $<
\pi$, ce que l'énoncé du lemme VIII ne donne pas. Il marque « faux » la
conclusion $]s', s''[\, \subset \operatorname{int}\Gamma$, et « ambigu » la
condition qui la précède. Page 36, au crayon, en face de « l'enveloppe convexe
$C$ de $(\mathring{T} \cup\, ]s', s''[) \cap \Gamma$ » : « ce n'est pas un
“ens.\ convexe” ! » ; en effet cette partie de $\Gamma$ est infinie, et la
convention du § 2 réserve le mot aux enveloppes d'ensembles finis, de sorte que
les lemmes IV et V ne s'appliquent pas tels quels\footnote{Les marques au
crayon des pages 35 et 36 sont d'une main que la transcription ne distingue pas
avec certitude de celle de l'encre noire.}.

Sa page 29, au crayon, au verso du feuillet de la page 28, esquisse la bonne
version. Soit $s$ un sommet en lequel le secteur d'angle $> \pi$ découpé par
les deux arêtes est extérieur à $\Gamma$ — un sommet d'angle intérieur $< \pi$,
qui existe toujours —, $s', s''$ ses voisins, $T$ leur triangle, et supposons
$n \geq 4$.
\begin{itemize}
  \item[(i)] Si $\Gamma \cap T = [s, s'] \cup [s, s'']$, alors $]s', s''[\,
    \subset \operatorname{Int}\Gamma$, et $[s', s'']$ est une diagonale.
  \item[(ii)] Sinon, soit $\Phi$ l'ensemble des sommets de $\Gamma$ dans
    $\mathring{T}$ ; s'il n'est pas vide et si $t \in \Phi$ est à distance
    maximale de la droite $s's''$, alors $]s, t[\, \subset
    \operatorname{Int}\Gamma$.
\end{itemize}
\emph{Cor Main.} Ou bien $]s', s''[\, \subset \operatorname{Int}\Gamma$, ou
bien il existe un sommet $t$, non voisin de $s$, tel que $]s, t[\, \subset
\operatorname{Int}\Gamma$ : tout polygone à au moins quatre côtés a une
diagonale intérieure\footnote{Sur la page, (i) écrit « $\Gamma \cap \mathring{T}
= [s, s'] \cup [s, s'']$ », le $\mathring{T}$ récrit sur un mot noirci ; c'est
$\Gamma \cap T$ qui est visé, comme dans (ii). Le Cor Main dit « $t \neq s$ » ;
il faut $t$ non voisin de $s$, puisque $]s, s'[$ est une arête. Enfin, il
manque un cas : si $\Phi = \emptyset$ mais que $\Gamma$ touche $]s', s''[$, on
prend pour $t$ un sommet de $\Gamma$ sur $]s', s''[$ (il y en a, une arête ne
pouvant traverser ni $[s, s']$ ni $[s, s'']$), et la conclusion tient. Le
complément est de cette lecture ; la page s'arrête au début de (ii), et le
passage barré du haut n'est lu qu'en partie.}.
La raison de (ii) est qu'aucune arête ne peut traverser $]s, t[$ sans avoir une
extrémité dans le petit triangle découpé dans $T$ par la parallèle à $s's''$
menée par $t$, où il n'y a pas de sommet. C'est le lemme classique qui sert à
trianguler les polygones, et dont on tire le théorème des deux
oreilles\footnote{Sous cette forme on le trouve par exemple chez N.~J.~Lennes
(1911) pour la triangulation, et dans la démonstration de Meisters ; références
de mémoire.}.

\pagerange{37}{47}
\subsection*{La récurrence, et le second théorème (pages 37 à 47)}

Le § 3 du mémoire, « Preuve du théorème 1 » (pages 37 à 43), prend $s, s'$
comme au corollaire de la page 36 — une diagonale extérieure, portée par
$\partial E(\Gamma)$ —, coupe $\Gamma$ en deux arcs $\Gamma^1_{ss'},
\Gamma^2_{ss'}$, et forme les polygones $\Lambda^1, \Lambda^2$ en leur
recollant $[s, s']$. Par récurrence sur $\operatorname{card} S(\Gamma) - 2$,
$\Lambda^1$ et $\Lambda^2$ sont des J-polygones ; les \emph{lemmes XII à XV}
établissent que l'un contient l'autre ($\operatorname{int}\Lambda^2 \subset
\operatorname{int}\Lambda^1$), et les \emph{lemmes XVI et XVII} construisent
les deux composantes de $\mathbb{R}^2 \smallsetminus \Gamma$ et vérifient la
condition locale. C'est le cas b) de la page 18. Ses marques : « $a = s$ ? $b =
s'$ » sur une notation flottante du lemme XII, et, au lemme XIII, la
simplification $\operatorname{Ext}(\partial E(\Gamma)) = \mathbb{R}^2
\smallsetminus E(\Gamma) \subset \operatorname{Ext}(\Lambda^i)$ — juste, et
sans disque : le complémentaire de $E(\Gamma)$ est connexe, non borné, disjoint
de $\Lambda^i$\footnote{Ces deux marques sont au crayon. Les pages 38, 39, 41,
42 et 44 ne portent aucune annotation en mots.}.

Le § 4, « Preuve du théorème 2 » (pages 43 à 47), coupe cette fois selon une
diagonale intérieure (lemme XI) ; les \emph{lemmes XVIII et XIX} établissent
que $\Lambda^1$ et $\Lambda^2$ sont extérieurs l'un à l'autre — le cas a) de la
page 18 —, le \emph{lemme XX} donne deux bons sommets à triangle intérieur, le
\emph{lemme XXI} un bon sommet à triangle extérieur pour un polygone non
convexe, et le théorème 2 en résulte\footnote{La note de la transcription sur
la page 43 glose le § 4 par « tout polygone est un J-polygone », qui est le
théorème 1 ; on suit la page 20 et l'en-tête de la transcription, pour lesquels
le § 4 démontre le théorème 2.}. Les pages 46 et 47 portent des annotations au
crayon d'une main que la transcription n'établit pas : au lemme XX,
l'insertion « (non consécutifs si $n \geq 4$ » — deux oreilles d'un polygone
à au moins quatre côtés peuvent en effet être choisies non adjacentes, ce qui
est la forme « sans chevauchement » du théorème de Meisters —, un « ? » sur un
$s''$ que le texte ne définit pas, « so what », « ou bien, », et une note en
biais lue en partie, « faux, si on … choisit … »\footnote{La page 45, verso de
la page 44, porte deux croquis au crayon de polygones non convexes, de sommets
$s', s, s''$, sans texte ; elle n'a pas de marqueur de page dans la
transcription.}.

Ce que la lecture de Grothendieck dégage du mémoire tient en peu de mots :
le lemme VIII n'a pas besoin de son hypothèse, le lemme XI est faux tel qu'il
est énoncé et vrai pour un sommet bien choisi et $n \geq 4$, et le sommet
choisi doit l'être par la démonstration, non par l'énoncé, d'un lemme
antérieur. Ses feuillets 29 et 31 en donnent les versions justes.

\end{document}
